
\section{Boolean reflection and bounded logical compilation}
\label{sec:boolean}

\subsection{Boolean terms inside iPRA}

Write \(0\) and \(1=S0\) for the corresponding terms.  A term \(b\) is
Boolean when
\[
       \mathsf{Bool}(b)\ \defeq\ (b=0)\vee(b=1).
\]
This is quantifier free.  On Boolean inputs define
\[
\begin{aligned}
 \mathsf{bnot}(b) &\defeq 1\mathbin{\dotminus} b,\\
 \mathsf{band}(b,c) &\defeq b\cdot c,\\
 \mathsf{bor}(b,c) &\defeq
   1\mathbin{\dotminus}
      \bigl((1\mathbin{\dotminus}b)
            (1\mathbin{\dotminus}c)\bigr),\\
 \mathsf{bimp}(b,c) &\defeq
   \mathsf{bor}(\mathsf{bnot}(b),c).
\end{aligned}
\]
All operations are primitive recursive.  Their truth tables are finite
derivations from the defining equations of addition, multiplication, and
truncated subtraction.  They should be proved once in the object calculus.

The existing equality test
\[
 \mathsf{eqb}(a,b)
   =1\mathbin{\dotminus}
      \bigl((a\mathbin{\dotminus}b)
            +(b\mathbin{\dotminus}a)\bigr)
\]
already has checked theorems
\[
  a=b\ \vdash\ \mathsf{eqb}(a,b)=1,
  \qquad
  \neg(a=b)\ \vdash\ \mathsf{eqb}(a,b)=0.
\]
The arithmetic library also derives decidable equality
\((a=b)\vee\neg(a=b)\).  Thus equality atoms have an internal Boolean
characteristic term without invoking classical logic as an ambient principle.

\subsection{Reflection for every fixed quantifier-free formula}

\begin{theorem}[Quantifier-free Boolean reflection]
For every fixed quantifier-free formula
\(\varphi(\vec x):\Fm\,\PTm\,n\), one can construct a term
\(\chi_\varphi(\vec x):\PTm\,n\) and object-iPRA derivations of
\[
\begin{aligned}
 &\mathsf{Bool}(\chi_\varphi),\\
 &\varphi\ \iffc\ \chi_\varphi=1,\\
 &\neg\varphi\ \iffc\ \chi_\varphi=0.
\end{aligned}
\]
The construction commutes with simultaneous term substitution.
\end{theorem}

\begin{proof}
Proceed by structural recursion on the evidence that \(\varphi\) is
quantifier free.

For an equality \(s=t\), take
\(\chi_\varphi=\mathsf{eqb}(s,t)\).  Split by the internally derived
\((s=t)\vee\neg(s=t)\).  In the first branch
\(\mathsf{eqb}=1\); in the second it is \(0\).  Successor nonzeroness excludes
\(0=1\), so both directions of the two equivalences follow
intuitionistically.

For falsity take \(\chi_\bot=0\).  The formula \(0=1\) is contradictory, while
\(0=0\) is reflexive.

For conjunction, disjunction, and implication define
\[
\begin{aligned}
 \chi_{\alpha\wedge\beta}
    &=\mathsf{band}(\chi_\alpha,\chi_\beta),\\
 \chi_{\alpha\vee\beta}
    &=\mathsf{bor}(\chi_\alpha,\chi_\beta),\\
 \chi_{\alpha\to\beta}
    &=\mathsf{bimp}(\chi_\alpha,\chi_\beta).
\end{aligned}
\]
By the induction hypotheses both arguments are Boolean.  Perform the finite
four-case analysis on their values.  In each case the required equation is a
defining-equation calculation and the logical equivalence is an
intuitionistic propositional derivation.  The implication case is legitimate
because its antecedent and consequent are decidable; no general law of
excluded middle is assumed.

Substitution compatibility follows in parallel.  Equality tests and Boolean
operations commute definitionally with substitution, while the recursive
formula compiler follows the already proved substitution laws for
\(\Fm\).  Therefore
\[
    \chi_{\varphi[\sigma]}
      = \chi_\varphi[\sigma]
\]
as terms, or at worst by a short congruence derivation if the implementation
uses derived symbols.
\end{proof}

\begin{remark}
The theorem is a \emph{metalevel compiler indexed by a fixed formula}.  It
does not require a primitive-recursive universal evaluator for formulas.  Lean
recurses over the concrete syntax and returns a concrete term and concrete
proofs.
\end{remark}

\subsection{Bounded universal folds}

Let \(\varphi(i,\vec x)\) be reflected by
\(\chi(i,\vec x)\).  Define by primitive recursion
\[
\begin{aligned}
  A(0,\vec x)&=1,\\
  A(N+1,\vec x)&=A(N,\vec x)\cdot \chi(N,\vec x).
\end{aligned}
\]
For Boolean factors, \(A\) is Boolean.

\begin{theorem}[Bounded universal reflection]
iPRA proves
\[
 A(N,\vec x)=1
 \quad\iffc\quad
 \forall i\,(i<N\to\varphi(i,\vec x)).
\]
It also proves that \(A(N,\vec x)=0\) yields a least counterexample
\(i<N\) when one exists.
\end{theorem}

\begin{proof}
For the forward direction, fix \(i<N\).  A prefix-factor lemma states that if
a finite Boolean product is \(1\), every factor below its length is \(1\).
Prove that lemma by quantifier-free induction on the product length, using the
equation \(A(N+1)=A(N)\chi(N)\) and the finite cases \(i<N\) or \(i=N\).
Reflection then gives \(\varphi(i,\vec x)\).

For the reverse direction, work under the universal hypothesis.  Induct on
\(N\) with the quantifier-free invariant \(A(N,\vec x)=1\).  The base equation
is immediate.  At the successor step, instantiate the universal hypothesis at
\(N\), obtain \(\chi(N,\vec x)=1\) by reflection, and use the recursion
equation.

For a counterexample, search below \(N\) for a zero of \(\chi\).  The checked
least-search theorem returns a least index and proves its minimality.  If the
product is zero, a finite-product zero lemma proves that such an index exists;
alternatively the product itself may be implemented through an early-stopping
search, making existence immediate.
\end{proof}

The crucial point is that the induction predicate is the single equation
\(A(N,\vec x)=1\).  The quantified formula occurs only as a context hypothesis
or final conclusion.

\subsection{Bounded existential reflection and witness extraction}

Define
\[
   f(i,\vec x)=\mathsf{bnot}(\chi(i,\vec x))
\]
and let
\[
   \mu_\varphi(N,\vec x)
\]
be the least \(i<N\) with \(f(i,\vec x)=0\), or \(N\) if none exists.  This is
exactly the convention of the checked \code{Search.mu} construction.

\begin{theorem}[Bounded least witness]
iPRA proves:
\[
\begin{aligned}
 \exists i<N\,\varphi(i,\vec x)
   &\quad\iffc\quad \mu_\varphi(N,\vec x)<N,\\
 \mu_\varphi(N,\vec x)<N
   &\quad\imp\quad
     \varphi(\mu_\varphi(N,\vec x),\vec x),\\
 j<\mu_\varphi(N,\vec x)
   &\quad\imp\quad
     \neg\varphi(j,\vec x).
\end{aligned}
\]
\end{theorem}

\begin{proof}
If a witness \(i<N\) exists, reflection gives \(\chi(i,\vec x)=1\), hence
\(f(i,\vec x)=0\).  Apply \code{search_found}; it returns a least zero below
\(N\).  Conversely, \code{search_spec} says that a returned index below \(N\)
is a zero of \(f\), which for a Boolean \(\chi\) means
\(\chi=1\), and reflection yields \(\varphi\).  Minimality is
\code{search_minimal} followed by the same Boolean equivalence.
\end{proof}

This theorem is stronger than a bare bounded-existence equivalence: it gives a
canonical least witness without any choice principle.

\subsection{Fixed bounded alternation}

A formula with a fixed number of bounded quantifier alternations can be
compiled recursively.  Existential levels return a least successful move;
universal levels return either success of every branch or a least failing
challenge.  The result is a pair of finite certificates:

\begin{itemize}
\item a success strategy, containing a response to every bounded universal
challenge encountered after an existential move;
\item a failure counterstrategy, containing a challenge defeating every
candidate move.
\end{itemize}

Both are finite arrays with explicit bounds.  Their verification is a Boolean
finite fold.  This gives a constructive determinacy theorem for each fixed
bounded game; it does not assert determinacy of arbitrary infinite games.

\subsection{Recommended implementation boundary}

The Boolean compiler should be implemented over the actual
\code{Fm PTm n} syntax with an input proof \code{phi.QF}.  Its output should be a
structure:
\[
\begin{split}
\mathsf{ReflectedQF}(\varphi)=\{&
  \mathsf{bit}:\PTm\,n,\quad
  \mathsf{isBool}:\ []\vdash\mathsf{Bool}(\mathsf{bit}),\\
 &\mathsf{trueIff}:\ []\vdash
       \varphi\leftrightarrow(\mathsf{bit}=1),\\
 &\mathsf{falseIff}:\ []\vdash
       \neg\varphi\leftrightarrow(\mathsf{bit}=0)\}.
\end{split}
\]
The compiler is metalevel structural recursion.  It should not add a
\(\mathsf{reflect}\) axiom or an opaque semantic oracle.
